%! TEX root = ../m3-20-21.tex

\chapter{Transformata Z}

Această transformată se definește pe un caz discret, pornind de la:
\begin{definition}\label{def:semnal}\index{semnal!discret}
    Se numește \emph{semnal discret} o funcție $ x : \ZZ \to \CC $ dată
    de $ n \mapsto x_n $ sau, echivalent, $ x(n) $ ori $ x[n] $.
\end{definition}

Mulțimea semnalelor discrete se va nota cu $ S_d $, iar cele cu suport pozitiv
(nule pentru $ n < 0 $ se va nota $ S_d^+ $.

Un semnal particular este \emph{impulsul unitar discret} la momentul $ k $,
definit prin:\index{semnal!unitar}
\[
    \delta_k(n) = \begin{cases} 1, & n = k \\ 0 & n \neq k \end{cases},
\]
definit pentru $ k \in \ZZ $ fixat.

Pentru $ k = 0 $, vom nota $ \delta_0 = \delta $.

\begin{definition}\label{def:intirziat} \index{semnal!intîrziat}
    Fie $ x \in S_d $ și $ k \in \ZZ $ fixat. Semnalul $ y = (x_{n-k}) $ se numește
    \emph{întîrziatul lui $ x $ cu $ k $ momente}.
\end{definition}

O operație foarte importantă, pe care se vor baza unele proprietăți esențiale ale
transformatei $ Z $ este \emph{convoluția:} \index{convoluție}
\begin{definition}\label{def:convolutie}
    Fie $ x, y \in S_d $. Dacă seria $ \ds\sum_{k \in \ZZ} x_{n-k} y_k $
    este convergentă pentru orice $ n \in \ZZ $ și are suma $ z_n $, atunci semnalul
    $ z = (z_n) $ se numește \emph{convoluția} semnalelor $ x $ și $ y $ și
    se notează $ z = x \ast y $.
\end{definition}

Trei proprietăți imediate sînt:
\begin{itemize}
    \item $ x \ast y = y \ast x $;
    \item $ x \ast \delta = x $;
    \item $ (x \ast \delta_k)(n) = x_{n - k} $.
\end{itemize}

Ajungem acum la definiția principală:
\begin{definition}\label{def:z}\index{transformata!Z}
    Fie $ s \in S_d $, cu $ s = (a_n)_n $. Se numește \emph{transformata Z} sau
    \emph{transformata Laplace discretă} a acestui semnal funcția definită prin:
    \[
        L_s(z) = Z_s(z) = \sum_{n \in \ZZ} a_n z^{-n},
    \]
    care se definește în domeniul de convergență al seriei Laurent din definiție.
\end{definition}

Principalele proprietăți pe care le vom folosi în calcule sînt:
\begin{enumerate}[(1)]
    \item \textbf{Inversarea transformării $ Z $:} Fie $ s \in S_d^+ $, cu $ s = (a_n) $.
        Presupunem că $ Z_s(z) $ este olomorfă în domeniul $ |z| \in (r, R) $. Atunci
        putem recupera semnalul $ a_n $ prin formula:
        \[
            a_n = \dfrac{1}{2 \pi i} \int_{\gamma} z^{n-1} Z_s(z) dz, \quad n \in \ZZ,
        \]
        unde $ \gamma $ este discul de rază $ \rho \in (r, R) $.
    \item \textbf{Teorema de convoluție:} Fie $ s, t \in S_d^+ $. Atunci $ s \ast t = S_d^+ $
        și are loc $ L_{s \ast t} = L_s \cdot L_t $. În particular:
        \[
            L_{s \ast \delta_k}(z) = z^{-k} L_s(z), \quad k \in \ZZ;
        \]
    \item \textbf{Prima teoremă de întîrziere:} Pentru $ n\in \NN^\ast $:
        \[
            L_s( z_{t - n}) = z^{-n} L_s(f);
        \]
    \item \textbf{A doua teoremă de întîrziere (teorema de deplasare):}
        \[
            L_s(z_{t + n}) = z^n \cdot \Big( L_s(z) - \sum_{t = 0}^{n-1} z_t z^{-t} \Big), \quad n \in \NN^\ast.
        \]
\end{enumerate}

Cîteva transformate uzuale sînt:
\begin{tabular}{c | c}
    $ s $ & $ L_s $ \\
    \hline\hline
    $ \begin{cases}
        h_n = 0, & n < 0 \\
        h_n = 1, & n \geq 0
    \end{cases} $ & $ \dfrac{z}{z - 1} $ \\
          & \\
    $ \delta_k, k \in \ZZ $ & $ \dfrac{1}{z^k} $ \\
                            &\\
    $ s = (n)_{n \in \NN} $ & $ \dfrac{z}{(z - 1)^2} $ \\
                            &\\
    $ s = (n^2)_{n \in \NN} $ & $ \dfrac{z(z+1)}{(z - 1)^3} $ \\
                              &\\
    $ s = (a^n)_{n \in \NN}, a \in \CC $ & $ \dfrac{z}{z - a} $ \\
                                         &\\
    $ s = (e^{an})_{n \in \NN}, a \in \RR $ & $ \dfrac{z}{z - e^a} $ \\
                                            &\\
    $ s = (\sin(\omega n))_{n \in \NN}, \omega \in \RR $ & $ \dfrac{z \sin \omega}{z^2 - 2z \cos \omega + 1} $ \\
                                                         &\\
    $ s = (\cos(\omega n))_{n \in \NN}, \omega \in \RR $ & $ \dfrac{z(z - \cos \omega)}{z^2 - 2z \cos \omega + 1} $
\end{tabular}

%%%%%%%%%%%%%%%%%%%%%%%%%%%%%%%%%%%%%%%%%%%%%%%%%%%%%%%%%%%%%%%%%%%%%% 

\section{Exerciții}

\indent\indent 1. Să se determine semnalul $x \in S_d^+$, a cărui transformată $Z$ este 
dată de:
\begin{enumerate}[(a)]
    \item $\Lap_s (z) = \dfrac{z}{(z-3)^2}$ ;
    \item $\Lap_s (z) = \dfrac{z}{(z-1)(z^2 + 1)}$ ;
    \item $\Lap_s (z) = \dfrac{z}{(z-1)^2(z^2 + z - 6)}$ ;
    \item $\Lap_s (z) = \dfrac{z}{z^2 + 2az + 2a^2}, a > 0$ parametru.
\end{enumerate}

\emph{Soluție:}

(a) Avem:
\begin{align*}
    x_n &= \dfrac{1}{2\pi i} \int_{|z| = \rho} z^{n-1} \Lap_s (z) dz \\
        &= \Rez(z^{n-1} \Lap_s(z), 3) \\
        &= \Rez \Big( \dfrac{z^n}{(z-3)^2}, 3 \Big) \\
        &= \lim_{z \to 3} \Big( (z-3)^2 \cdot \dfrac{z^n}{(z-3)^2} \Big)^{\prime} \\
        &= \lim_{z \to 3} nz^{n-1} \\
        &= n3^{n-1}.
\end{align*}

(b)
\begin{align*}
    x_n &= \dfrac{1}{2 \pi i} \int_{|z| = \rho} z^{n-1} \Lap_s(z) dz \\
        &= \Rez(z^{n-1} \Lap_s (z), 1) + \Rez(z^{n-1} \Lap_s(z), i) + 
        \Rez(z^{n-1} \Lap_s(z), -i) \\
    \Rez(z^{n-1}\Lap_s(z),1) &= \lim_{z \to 1} z^{n-1} \dfrac{z}{(z - 1)(z^2 + 1)} 
    \cdot (z-1) = \dfrac{1}{2} \\
    \Rez(z^{n-1} \Lap_s(z), i) &= \dfrac{i^n}{2i \cdot (i-1)} \\
    \Rez(z^{n-1} \Lap_s(z), -i) &= \dfrac{(-1)^n i^n}{2i(i+1)}.
\end{align*}

Remarcăm că pentru $n = 4k$ și $n = 4k+1$, avem $x_n = 0$, iar în celelalte 
două cazuri, $x_n = 1$. \\

(c)
\begin{align*}
    \Rez(z^{n-1} \Lap_s(z), 1) &= \lim_{z\to 1} \Big[ z^{n-1} \cdot (z-1)^2 \cdot %
    \dfrac{z^2}{(z-1)^2 \cdot (z^2 + z - 6)} \Big]^{\prime}\\
                               &= - \dfrac{4n + 3}{16}. \\
        \Rez(z^{n-1} \Lap_s(z), 2) &= \dfrac{2^n}{5} \\
        \Rez(z^{n-1} \Lap_s(z), -3) &= -\dfrac{(-3)^n}{80}.
\end{align*}

Obținem $x_n = - \dfrac{4n+3}{16} + \dfrac{2^n}{5} - \dfrac{(-3)^n}{80}$. \\

(d) $z_{1,2} = a (-1 \pm i)$ sînt poli simpli. Avem:
\begin{align*}
    x_n &= \Rez\Big( \dfrac{z^n}{(z^2 + 2a + 2a^2)}, z_1 \Big) + 
    \Rez\Big(  \dfrac{z^n}{(z^2 + 2a + 2a^2)}, z_2 \Big) \\
        &= \dfrac{a^n(-1 + i)^n}{2z_1 + 2a} + \dfrac{a^n(-1-i)^n}{2z_1 + 2a} \\
        &= -\dfrac{i}{2a}(z_1^n - z_2^n).
\end{align*}

Putem scrie trigonometric numerele $z_1$ și $z_2$:
\begin{align*}
    z_1 &= a(-1 + i) = a \sqrt{2} \Big( \cos \dfrac{3\pi}{4} + i \sin \dfrac{3 \pi}{4} \Big) \\
    z_2 &= a(-1 - i) = a \sqrt{2} \Big( \cos \dfrac{3 \pi}{4} - i \sin \dfrac{3 \pi}{4} \Big).
\end{align*}

Deci: $x_n = 2^{\frac{n}{2}} a^{n-1} \sin \dfrac{3 n \pi}{4}$.


%%%%%%%%%%%%%%%%%%%%%%%%%%%%%%%%%%%%%%%%%%%%%%%%%%%%%%%%%%%%%%%%%%%%%%%%%%%%%%%% 

\vspace{0.3cm}

2. Fie $x = (x_n) \in S_d^+$ și $y = (y_n)$, unde $y_n = x_0 + \dots + x_n$. 
Să se arate că $Y(z) = \dfrac{z}{z-1} X(z)$.

\emph{Soluție:}

Avem $Y(z) = \displaystyle\sum_{n=0}^\infty y_n z^{-n}.$ Dar: 
\[
    X(z) = \sum_{n=0}^\infty x_n z^{-n} \text{ și } 
    \sum_{n=0}^\infty x_{n-1} z^n = \dfrac{1}{z} X(z),
\]
deoarece $x_{-1} = 0$. Putem continua și obținem 
$\displaystyle\sum_{n=0}^\infty x_{n-k} z^{-n} = \dfrac{1}{z^k} X(z)$. Așadar: 
\[
    Y(z) = X(z) \cdot \Big( 1 + \dfrac{1}{z} + \dfrac{1}{z^2} + \dots \Big) 
    = X(z) \cdot \dfrac{z}{z-1}.
\]


\vspace{0.3cm}

3. Cu ajutorul transformării $Z$, să se determine șirurile $(x_n)$ definite 
prin următoarele relații:
\begin{enumerate}[(a)]
    \item $x_0 = 0, x_1 = 1, x_{n+2} = x_{n+1} + x_n, n \in \mathbb{N}$ 
        (șirul lui Fibonacci);
    \item $x_0 = 0, x_1 = 1, x_{n+2} = x_{n+1} - x_n, n \in \mathbb{N}$;
    \item $x_0 = x_1 = 0, x_2 = -1, x_3 = 0, x_{n+4} + 
        2x_{n+3} + 3 x_{n+2} + 2 x_{n+1} + x_n = 0, n \in \mathbb{N}$;
    \item $x_0 = 2, x_{n+1} + 3x_n = 1, n \in \mathbb{Z}$;
    \item $x_0 = 0, x_1 = 1, x_{n+2} - 4 x_{n+1} + 3x_n 
        = (n+1)4^n, n \in \mathbb{N}$.
\end{enumerate}

\emph{Soluție:}

Abordarea generală este să considerăm șirul $(x_n)$ ca fiind restricția unui 
semnal $x \in S_d^+$ la $\mathbb{N}$ și rescriem relațiile de recurență sub 
forma unor ecuații de convoluție $a \ast x = y$, pe care le rezolvăm în $S_d^+$.\\

(a) Fie $x \in S_d^+$, astfel încît restricția lui la $\mathbb{N}$ să fie șirul 
căutat. Deoarece avem: 
\[
    x_{n+2} - x_{n+1} - x_n = y_n, n \in \mathbb{Z},
\]
cu $y_n = 0$ pentru $n \neq -1$ și $y_{-1} = 1$, avem ecuația de convoluție: 
\[
    a \ast x = y, \text{ unde } a 
    = \delta_{-2} + \delta_{-1} + \delta, y = \delta_{-1}.
\]

Aplicăm transformata $Z$ și rezultă: 
\[
    \Lap_s x(z) (z^2 - z - 1) = z \Longrightarrow x_n 
    = \dfrac{1}{\sqrt{5}} \Big[ \Big( \dfrac{1 + \sqrt{5}}{2} \Big)^n - 
    \Big( \dfrac{1 - \sqrt{5}}{2} \Big)^n \Big].
\]

(b) Ca în cazul anterior, avem $a \ast x = y$, cu 
$a = \delta_{-2} - \delta_{-1} + \delta$, unde $y = \delta_{-1}$. Aplicînd 
transformata $Z$, obținem: 
\[
    \Lap_s x(z) (z^2 - z + 1) = z \Longrightarrow \Lap_s x(z) 
    = \dfrac{z}{z^2 - z + 1}.
\]

Obținem:
\[
    x_n = \Rez \Big( z^{n-1} \dfrac{z}{z^2 - z + 1}, 
        \dfrac{1 + i \sqrt{3}}{2} \Big) + \Rez \Big( z^{n-1} \dfrac{z}{z^2 - z + 1}, 
    \dfrac{1 - i \sqrt{3}}{2} \Big).
\]

Calculăm reziduurile, cu notația $\varepsilon = \dfrac{1 + i \sqrt{3}}{2}$ și 
$\overline{\varepsilon} = \dfrac{1 - i \sqrt{3}}{2}$:
\begin{align*}
    \Rez \Big( \dfrac{z^n}{z^2 - z + 1}, \varepsilon \Big) 
   &= \lim_{z \to \varepsilon} \dfrac{z^n}{z^2 - z + 1} 
   \big( z - \varepsilon \big) \\
   &= \dfrac{\varepsilon^n}{i \sqrt{3}} \\
   &= \dfrac{\cos \frac{2 n \pi}{3} + i \sin \frac{2 n \pi}{3}}{i \sqrt{3}}.
\end{align*}

Similar: 
\[
    \Rez \Big( \dfrac{z^n}{z^2 - z + 1}, \overline{\varepsilon} \Big) 
    = \dfrac{\cos \frac{2n\pi}{3} - i \sin \frac{2 n \pi}{3}}{-i\sqrt{3}}.
\]

Rezultă: 
\[
    x_n = \dfrac{2}{\sqrt{3}} \sin \dfrac{2n\pi}{3}, n \in \mathbb{N}.
\]

(c) Ecuația $a \ast x = y$ este valabilă pentru: 
\[
    a = \delta_{-4} + 2 \delta_{-3} + 3 \delta_{-2} + 2 \delta_{-1} + \delta, 
    \quad y = -\delta_{-2} - 2 \delta_{-1}.
\]

Aplicăm transformata $Z$ și obținem: 
$\Lap_s x(z) = - \dfrac{z(z+2)}{(z^2 + z + 1)^2}$. Descompunem în fracții 
simple, calculăm reziduurile și ținem cont de faptul că rădăcinile numitorului, 
$\varepsilon_1, \varepsilon_2$ sînt poli de ordinul 2, obținem: 
\[
    x_n = \dfrac{(2n-4)(\varepsilon_1^n - \varepsilon_2^n) - 
        (n+1)(\varepsilon_1^{n-1} + \varepsilon_1^{n-2} - \varepsilon_2^{n-1} - 
    \varepsilon_2^{n-2})}{(\varepsilon_2 - \varepsilon_1)^3} 
    = \dfrac{2(n-1)}{\sqrt{3}} \sin \dfrac{2 n \pi}{3}, n \in \mathbb{N}.
\]

(d) Ecuația corespunzătoare este $a \ast x = y$, cu $a = \delta_{-1} + 3 \delta$ și 
$y_n = 1, \forall n \geq 1$, iar $y_{-1} = x_0 + 3 x_{-1} = 2$, cu 
$y_n = 0, \forall n \leq -2$, adică $y = 1 + 2 \delta_{-1}$.

Așadar: 
\[
    \delta_{-1} \ast x + 3 \delta \ast x = 1 + 2 \delta_{-1}.
\]

Aplicăm transformata $Z$ și obținem:
\begin{align*}
    z \Lap_s x(z) + 3 \Lap_s x(z) &= \dfrac{z}{z-1} + 2z \\
                                  &= \dfrac{2z^2 + 3z}{z-1} \\
    \Longrightarrow \Lap_s x(z) &= \dfrac{2z^2 + 3z}{(z-1)(z+3)} \\
    \Longrightarrow x_n &= \Rez(z^{n-1} \cdot \dfrac{2z^2 + 3z}{(z-1)(z+3)}, 1) + 
    \Rez (z^{n-1} \cdot \dfrac{2z^2 + 3z}{(z-1)(z+3)}, -3) \\
    \Rez(z^{n-1} \cdot \dfrac{2z^2 + 3z}{(z-1)(z+3)}, 1) &= \lim_{z \to 1} (z-1) 
    \cdot z^{n-1} \cdot \dfrac{2z^2 + 3z}{(z-1)(z+3)} = \dfrac{5}{4} \\
    \Rez(z^{n-1} \cdot \dfrac{2z^2 + 3z}{(z-1)(z+3)}, -3) 
                                                         &= \lim_{z \to -3} (z+3) z^{n-1} \cdot \dfrac{2z^2 + 3z}{(z-1)(z+3)} \\
                                                         &= (-3)^{n-1} \cdot \dfrac{12}{-4} = -3 \cdot (-3)^{n-1}.	                              
\end{align*}

Rezultă: $x_n = \dfrac{5}{4} - 3 \cdot (-3)^{n-1}$. \\

(e) Avem ecuația: $a \ast x = y$, unde 
$a = \delta_{-2} - 4 \delta_{-1} + 3 \delta$, cu 
$y_n = 0, \forall n \leq -2, y_{-1} = 1$ și 
$y_n = (n+1)4^n, \forall n \in \mathbb{N}$.

Fie $s_1 = (n4^n)_n, s_2 = (4^n)_n$. Atunci:
\begin{align*}
    \Lap_s s_1(z) &= -z \Lap_s s_2^{\prime}(z) 
    = - z \big( \dfrac{z}{z-4} \big)^{\prime} = \dfrac{4z}{(z-4)^2} \\
    \Longrightarrow \Lap_s x(z)(z^2 - 4z + 3) 
                  &= \dfrac{4z}{(z-4)^2} + \dfrac{z}{z-4} + z \\
                  &= \dfrac{z^2}{(z-4)^2} + z \\
    \Longrightarrow  \Lap_s x(z) &= \dfrac{z(z^2 - 7z + 16)}{(z-4)^2 (z-1) (z-3)}.
\end{align*}

Descompunem în fracții simple și obținem, în fine: 
\[
    x_n = \dfrac{1}{9} \big[ 18 \cdot 3^n + (3n - 13) 4^n - 5 \big], 
    n \in \mathbb{N}.
\]

\textbf{OBSERVAȚIE:} Toate exercițiile cu recurențe se mai pot rezolva în alte
două moduri:

(1) Se poate aplica teorema de convoluție relației de recurență. De exemplu,
din recurența:
\[
    x_{n + 2} - 2x_{n + 1} + x_n = 2
\]
putem obține:
\[
    Z(x_{n+2}) - 2Z(x_{n + 1}) + Z(x_n) = 2Z(1),
\]
iar $ Z(x_{n+2}) = Z(x_n \ast \delta_{-2}) = Z(x_n) \cdot Z(\delta_{-2}) $ etc.

(2) Se poate aplica teorema de deplasare. În aceeași recurență de mai sus,
de exemplu, avem:
\[
    Z(x_{n+2}) = z^n \Big( Z(x_n) - x_0 - x_1 z^{-1} \Big)
\]
și la fel pentru celelalte.
